\begin{frame}{각 스텝은 ``앞 단계의 잔차만'' 배운다}
  \begin{columns}[T]
    \begin{column}{0.52\textwidth}
      \centering
      \includegraphics[width=\linewidth]{ch07_gbdt_residual_fitting.png}
      \vskip0.3em
      {\footnotesize $y=x$를 스템 1·2·4·8개의 합으로 근사 —
      각 단계의 계단형 합이 직선에 가까워진다.}
    \end{column}
    \begin{column}{0.45\textwidth}
      \begin{itemize}
        \item 스템(2단계 함수)은 \textbf{표현력이 거의 없는 기초 블릭} —
          그 잔차를 \textbf{순서대로} 추격하는 것만으로 직선(심지어
          곡선)을 자유롭게 근사
        \item 랜덤 포레스트 = 여러 트리의 \textbf{다수결}
        \item GBDT = 여러 트리의 \textbf{연속 수정}(correction)
        \item \textbf{같은 트리 합}이지만, 그 합이 만들어지는
          \textbf{방식이 정반대}
      \end{itemize}
    \end{column}
  \end{columns}
\end{frame}
